%% hopfion-framework/quantum/hilbert_space.tex
%% Sources: main_paper10.tex
%% Condensate Hilbert space, geometric quantisation, Born rule, golden-ratio spectrum
%% Auto-generated from project files. Edit source papers to update.

%════════════════════════════════════════════════════════════════════
% Geometric Quantisation — Paper X
%════════════════════════════════════════════════════════════════════
\begin{theorem}[Condensate symplectic structure]
\label{P10:thm:symplectic}
The condensate configuration space $\CQ$ carries the closed 2-form
\begin{equation}
  \bOmega
  \;=\; \int_{\RR^3} \delta\pi(x)\wedge\delta\rho(x)\,d^3x
  \;\approx\; \int_{\RR^3} W(x)\,\delta\dot\rho(x)\wedge\delta\rho(x)\,d^3x,
  \label{P10:eq:symplectic}
\end{equation}
where $W(x) = 2/(\vp^6(1+\bst|\nabla f_*|^2)^2)>0$ pointwise.
The kernel of $\bOmega$ at $f_*$ is the one-dimensional
subspace spanned by the Derrick scale mode:
\begin{equation}
  \ker\bOmega\big|_{f_*}
  \;=\; \mathrm{span}\bigl\{g\bigr\},
  \qquad
  g \;=\; -(r\partial_r f_* + z\partial_z f_*).
  \label{P10:eq:kernel}
\end{equation}
\end{theorem}
\status{published}
\source{P10:thm:symplectic}

\begin{remark}[Connection to Paper~IX]
\label{P10:rem:connection_paper9}
The force projection $F\mapsto F - \langle F,g\rangle/\langle g,g\rangle\cdot g$
of Paper~IX~\cite{Paper9} is precisely the projection onto
$(\ker\bOmega)^\perp$.
Theorem~\ref{P10:thm:symplectic} gives this an intrinsic geometric
interpretation: the projection removes the degenerate direction of
$\bOmega$ and restricts the gradient flow to the symplectic
leaf through $f_*$.
\end{remark}
\status{published}
\source{P10:rem:connection_paper9}

\begin{definition}[Derrick dilation group]
\label{P10:def:derrick_dilation}
Let $\mathrm{Diff}_1(\RR^3)$ denote the one-parameter group of
radial dilations $\sigma_a: f(r,z)\mapsto f(r/a,z/a)$, $a>0$.
Its Lie algebra is generated by $g = -(r\partial_r + z\partial_z)$.
\end{definition}
\status{published}
\source{P10:def:derrick_dilation}

\begin{theorem}[Reduced phase space]
\label{P10:thm:reduction}
The quotient
\begin{equation}
  \CQr \;\equiv\; \CQ\,/\,\mathrm{Diff}_1(\RR^3)
  \label{P10:eq:reduced}
\end{equation}
is a Banach manifold.
The reduced symplectic form
$\bOmega^{\mathrm{red}} = i^*\bOmega$ (pullback
to the slice $\{f\in\CQ : \sopt[f] = 2^{1/6}\}$)
is non-degenerate.
The condensate saddle $f_*$ is a fixed point of the
$\mathrm{Diff}_1$-action and projects to a point
$[f_*]\in\CQr$.
\end{theorem}
\status{published}
\source{P10:thm:reduction}

\begin{remark}
\label{P10:rem:phase_moment_map}
The phase $\Phi = Q\cdot\omega$ of Paper~IX~\cite{Paper9} is the
value of the moment map $\mu$ at the saddle configuration.
The \textsc{u}(1) symmetry generated by $\mu$ is the symmetry whose
unique invariant measure is used to derive the Born rule in
Theorem~\ref{P10:thm:quantisation} below.
\end{remark}
\status{published}
\source{P10:rem:phase_moment_map}

\begin{theorem}[Prequantum line bundle]
\label{P10:thm:prequantum}
The reduced symplectic manifold $(\CQr,\bOmega^{\mathrm{red}})$
admits a prequantum Hermitian line bundle
$\mathcal{L}\to\CQr$ with connection $\nabla$ of curvature
$F_\nabla = -i\bOmega^{\mathrm{red}}$
if and only if
\begin{equation}
  \left[\frac{\bOmega^{\mathrm{red}}}{2\pi}\right]
  \;\in\; H^2(\CQr,\ZZ).
  \label{P10:eq:weil}
\end{equation}
This condition is satisfied because $Q\in\ZZ$.
\end{theorem}
\status{published}
\source{P10:thm:prequantum}

\begin{remark}[The Hopf charge as quantisation condition]
\label{P10:rem:hopf_charge_quantisation}
Equation~\eqref{P10:eq:integrality} shows that the integrality of the
Hopf charge $Q\in\ZZ=\pi_3(S^2)$ is \emph{not} an independent
assumption: it is equivalent to the Weil integrality condition
that allows the condensate to be quantised.
A condensate with non-integer $Q$ (topologically impossible) could
not be consistently quantised.
This provides a \emph{topological reason} for quantisation:
the discreteness of $\pi_3(S^2)$ forces the discreteness of the
quantum spectrum.
\end{remark}
\status{published}
\source{P10:rem:hopf_charge_quantisation}

\begin{theorem}[Quantisation and Born rule]
\label{P10:thm:quantisation}
Let $(\CQr,\bOmega^{\mathrm{red}},\mathcal{L})$ be the
prequantum data of Theorem~\ref{P10:thm:prequantum}.
Choose as polarisation the real polarisation
$\mathcal{P} = T^{\mathrm{vert}}\CQr$ (vertical tangent bundle
of the Lagrangian fibration defined by the moment map $\mu$).
Then:

\begin{enumerate}[label=\emph{(\roman*)}]
  \item \emph{(Hilbert space.)}
    The geometric quantisation of $(\CQr,\bOmega^{\mathrm{red}},
    \mathcal{L},\mathcal{P})$ produces the Hilbert space
    \begin{equation}
      \Hil \;=\; L^2(\CQr,\dmu),
      \label{P10:eq:Hilbert}
    \end{equation}
    where $\dmu = \bOmega^{\mathrm{red}\,\wedge n}/n!$
    is the Liouville measure and $n = \tfrac{1}{2}\dim\CQr$.

  \item \emph{(Born rule.)}
    The inner product on $\Hil$ is the unique \textsc{u}(1)-invariant
    inner product:
    \begin{equation}
      \langle\psi,\phi\rangle
      \;=\; \int_{\CQr} \overline{\psi([f])}\,\phi([f])\,\dmu([f]),
      \label{P10:eq:inner_product}
    \end{equation}
    and the probability of outcome $a$ given state $\psi$ is
    \begin{equation}
      P(a\,|\,\psi) \;=\; |\langle e_a,\psi\rangle|^2,
      \label{P10:eq:born}
    \end{equation}
    where $\{e_a\}$ is an orthonormal basis of $\Hil$.
    This is the Born rule $P=|\psi|^2$.
\end{enumerate}
\end{theorem}
\status{published}
\source{P10:thm:quantisation}

\begin{remark}[Why $Q=10$ and not another integer]
\label{P10:rem:why_Q10}
The theorem holds for any $Q\in\ZZ\setminus\{0\}$.
The specific value $Q=10$ is determined by the WZW conformal
field theory of the condensate~\cite{Paper1}: the Hopfion carries
charge $Q = 2\cdot\text{rank}(G_2) + \text{(generation sum)} = 10$
where $G_2$ is the automorphism group of the octonions~\cite{Paper9}.
The quantisation is thus not arbitrary; $Q=10$ is forced by the
same algebraic structure that fixes the three lepton generations.
\end{remark}
\status{published}
\source{P10:rem:why_Q10}

\begin{theorem}[Spectrum of $\hat{H}_\mathrm{fb}$]
\label{P10:thm:spectrum}
The eigenvalues of $\hat{H}_\mathrm{fb}$ on $\Hil = L^2(\CQr,\dmu)$
form the discrete set
\begin{equation}
  \boxed{\mathrm{spec}(\hat{H}_\mathrm{fb})
  = \bigl\{E_n = \vp^{2n}E_0 : n\in\ZZ_{\geq0}\bigr\}},
  \label{P10:eq:spectrum}
\end{equation}
where $E_0 = \Efb[f_*;\bst]$ is the saddle energy and
$\vp^2 = \vp^2$ is the golden-ratio squared.
\end{theorem}
\status{published}
\source{P10:thm:spectrum}

\begin{remark}[Discreteness from topology]
\label{P10:rem:discreteness_from_topology}
The discreteness of the spectrum~\eqref{P10:eq:spectrum} follows ultimately
from the discreteness of $Q\in\ZZ=\pi_3(S^2)$.
The continuous limit $Q\to$ real number would give a continuous
spectrum, but $\pi_3(S^2)=\ZZ$ forbids non-integer $Q$.
Discreteness of the Hopf charge thus implies discreteness of the
quantum energy spectrum: the quantisation of charge implies
the quantisation of energy.
\end{remark}
\status{published}
\source{P10:rem:discreteness_from_topology}

\begin{theorem}[Resolution of Conjecture~1 of Paper~IX]
\label{P10:thm:main}
The canonical (geometric) quantisation of the density-feedback
Hopfion condensate with action~\eqref{P10:eq:L_full} produces:
\begin{enumerate}[label=\emph{(\roman*)}]
  \item The Hilbert space $\Hil = L^2(\CQr,\dmu)$,
    where $\CQr$ is the reduced configuration space
    (configuration space modulo Derrick dilations)
    and $\dmu$ is the Liouville measure of the reduced
    symplectic form (Theorem~\ref{P10:thm:quantisation}).

  \item The Born rule $P=|\psi|^2$ as the unique
    \textsc{u}(1)-invariant probability measure
    on $\Hil$ (Theorem~\ref{P10:thm:quantisation}).

  \item A discrete energy spectrum $E_n = \vp^{2n}E_0$
    consistent with the golden-ratio tower of Papers~IV--VI
    (Theorem~\ref{P10:thm:spectrum}).

  \item A $q$-deformed oscillator algebra with $q=\vp^2$
    on $\Hil$ (Theorem~\ref{P10:thm:q_algebra}), with
    $q$-numbers $[n]_{\vp^2}=F_{2n}$ (even Fibonacci numbers),
    giving closed-form scattering amplitudes.

  \item Zero one-loop correction to the Bekenstein--Hawking entropy:
    $\delta S_\mathrm{1-loop}=0$
    (Theorem~\ref{P10:thm:entropy_correction}), together with a
    logarithmically soft heat kernel
    $\Theta(t)\sim-\log(t)/(2\log\vp)$
    (Proposition~\ref{P10:prop:heat_kernel}).

  \item Exact Bekenstein--Hawking entropy at tree level:
    $S_\mathrm{Wald}=A_H/(4G_N)$ (Theorem~\ref{P10:thm:wald}),
    proved via the Wald formula: the $k$-essence Lagrangian
    $\partial\mathcal{L}_\mathrm{cond}/\partial R_{abcd}=0$
    contributes zero, leaving only the Einstein--Hilbert term.
\end{enumerate}
The \textsc{u}(1) symmetry group is generated by the moment map
$\mu = Q\cdot\omega$, where $Q=10\in\ZZ=\pi_3(S^2)$ is the
Hopf charge.
The prequantum line bundle exists because $Q\in\ZZ$
(Theorem~\ref{P10:thm:prequantum}): the topological quantisation
of the Hopf charge is equivalent to the Weil integrality
condition for geometric quantisation.
\end{theorem}
\status{published}
\source{P10:thm:main}

\begin{theorem}[Self-adjointness and spectral gap]
\label{P10:thm:selfadjoint}
The operator $\mathcal{L}_*$ is self-adjoint on $H^2(\RR^3)$
with domain $D(\mathcal{L}_*) = H^2(\RR^3)\ominus\mathrm{span}\{g\}$.
Its spectrum consists entirely of discrete eigenvalues
$0 < \lambda_0 \leq \lambda_1 \leq \cdots \to \infty$,
with no essential spectrum.
The spectral gap satisfies
\begin{equation}
  \lambda_0 \;=\; E_0\,\vp^0 \;=\; E_0,
  \qquad
  \Delta_\mathrm{spec} \;\equiv\; \lambda_1 - \lambda_0
  \;=\; E_0(\vp^2-1) \;=\; E_0\,\vp,
  \label{P10:eq:spectral_gap}
\end{equation}
where $E_0 = \Efb[f_*;\bst]$ is the saddle energy and
$\vp^2 - 1 = \vp$ is the golden-ratio identity.
\end{theorem}
\status{published}
\source{P10:thm:selfadjoint}

\begin{remark}[The spectral gap is golden]
\label{P10:rem:spectral_gap_golden}
The spectral gap equals $E_0\vp$ --- the saddle energy multiplied by the
golden ratio itself.
This is a direct consequence of the WZW conformal structure:
the gap between the ground state and the first excited state is
exactly the ratio $\vp^2/\vp^0 - 1 = \vp^2 - 1 = \vp$.
No other positive real number has its square exceed it by exactly itself.
\end{remark}
\status{published}
\source{P10:rem:spectral_gap_golden}

\begin{theorem}[Spectral zeta function]
\label{P10:thm:spectral_zeta}
The spectral zeta function of $\mathcal{L}_*$ is
\begin{equation}
  \boxed{\zeta_{\mathcal{L}_*}(s)
  \;\equiv\; \sum_{k=0}^{\infty} \lambda_k^{-s}
  \;=\; \frac{E_0^{-s}}{1 - \vp^{-2s}}},
  \qquad \mathrm{Re}(s) > 0.
  \label{P10:eq:zeta_exact}
\end{equation}
This admits a meromorphic continuation to all $s\in\CC$ with a unique
simple pole at $s=0$ of residue $1/(2\log\vp)$.
The zeta-regularised functional determinant is
\begin{equation}
  \det(\mathcal{L}_*) \;=\; \exp\!\bigl(-\zeta'_{\mathcal{L}_*}(0)\bigr)
  \;=\; \vp^{-1/6}\,E_0^{1/2},
  \label{P10:eq:det_exact}
\end{equation}
where the prime denotes the derivative with respect to $s$ at $s=0$.
The Liouville measure normalisation factor is
\begin{equation}
  \det(\mathcal{L}_*)^{-1/2}
  \;=\; \vp^{1/12}\,E_0^{-1/4}.
  \label{P10:eq:det_half}
\end{equation}
\end{theorem}
\status{published}
\source{P10:thm:spectral_zeta}

\begin{theorem}[$q$-oscillator algebra of the quantised Hopfion]
\label{P10:thm:q_algebra}
Define ladder operators on $\Hil = \mathrm{span}\{|n\rangle : n\geq0\}$ by
\begin{equation}
  a\,|n\rangle = \sqrt{[n]_q}\,|n-1\rangle,
  \qquad
  a^\dagger|n\rangle = \sqrt{[n+1]_q}\,|n+1\rangle,
  \qquad
  N\,|n\rangle = n\,|n\rangle,
  \label{P10:eq:ladder_def}
\end{equation}
with $q = \vp^2$ and $a|0\rangle = 0$.
These operators satisfy:
\begin{enumerate}[label=(\roman*),noitemsep]
\item $H_\mathrm{fb} = E_0\,q^N$ \quad (the Hamiltonian in terms of $N$);
\item $[N,\,a] = -a$, \quad $[N,\,a^\dagger] = +a^\dagger$;
\item \textbf{Macfarlane--Biedenharn relation:}
  \begin{equation}
    \boxed{a\,a^\dagger \;-\; q\,a^\dagger a
    \;=\; q^{-N}};
    \label{P10:eq:qosc}
  \end{equation}
\item $H_\mathrm{fb}\,a = q^{-1}\,a\,H_\mathrm{fb}$,
  \quad $H_\mathrm{fb}\,a^\dagger = q\,a^\dagger H_\mathrm{fb}$.
\end{enumerate}
\end{theorem}
\status{published}
\source{P10:thm:q_algebra}

\begin{remark}[Connection to the SU(2)$_3$ WZW model]
\label{P10:rem:jones}
The $q$-deformation parameter $q = \vp^2$ of Theorem~\ref{P10:thm:q_algebra}
is the \emph{square} of the Jones polynomial parameter at the fifth root
of unity: $q_\mathrm{Jones} = e^{i\pi/5}$ satisfies
$q_\mathrm{Jones} + q_\mathrm{Jones}^{-1} = 2\cos(\pi/5) = \vp$,
so that the quantum dimension of the spin-$\tfrac{1}{2}$ representation
of $\mathrm{SU}(2)_3$ is $\vp$ --- the same value that appears as the
RG fixed point and the Verlinde channel count $N_{\tau\to\mu} = \vp$
of Paper~V~\cite{Paper5}.
The real $q$-oscillator parameter $q = \vp^2$ is the analytic continuation
of $q_\mathrm{Jones}$ to the positive real axis, consistent with the
Wick-rotated (Euclidean) quantisation of the condensate.
This provides an independent confirmation that the $\mathrm{SU}(2)_3$ WZW
level $k=3$ fixed by Papers~I--IV is the correct WZW structure for the
$q$-algebra of the quantised Hopfion.
\end{remark}
\status{published}
\source{P10:rem:jones}

\begin{theorem}[Vanishing of the one-loop bulk entropy correction]
\label{P10:thm:entropy_correction}
\begin{equation}
  \boxed{\delta S_\mathrm{1-loop}
  \;\equiv\; -\frac{\partial\,\Gamma_\mathrm{1-loop}}{\partial T_H}
  \;=\; 0}.
  \label{P10:eq:delta_S_zero}
\end{equation}
The Bekenstein--Hawking entropy $S_\mathrm{BH} = A/(4G_N)$
is reproduced at tree level from Paper~VII~\cite{Paper7},
with no quantum correction from bulk condensate modes at one loop.
\end{theorem}
\status{published}
\source{P10:thm:entropy_correction}

\begin{remark}[Entropy protected by the golden-ratio spectrum]
\label{P10:rem:entropy_protected}
Theorem~\ref{P10:thm:entropy_correction} shows that the golden-ratio
spectrum \emph{protects} $S = A/4$ against one-loop quantum corrections.
In the Wald entropy framework~\cite{Wald1993}, the complete tree-level
calculation is carried out in Theorem~\ref{P10:thm:wald}:
the $k$-essence sector contributes zero to the Wald entropy because
$\partial\mathcal{L}_\mathrm{cond}/\partial R_{abcd}=0$, and the
one-loop correction also vanishes.
The Bekenstein--Hawking entropy $S = A_H/(4G_N)$ is therefore
an \emph{exact} result of the condensate framework at both tree level
and one loop.
\end{remark}
\status{published}
\source{P10:rem:entropy_protected}

\begin{theorem}[Wald entropy of the acoustic condensate]
\label{P10:thm:wald}
For the total action~\eqref{P10:eq:action_total}, the Wald entropy of any
stationary black-hole solution is
\begin{equation}
  \boxed{S_\mathrm{Wald}
  \;=\; \frac{A_H}{4G_N}},
  \label{P10:eq:wald_result}
\end{equation}
where $A_H$ is the horizon area in the physical (Kerr) metric.
The condensate $k$-essence sector contributes exactly zero to the
Wald entropy at tree level.
\end{theorem}
\status{published}
\source{P10:thm:wald}

\begin{remark}[Acoustic metric and the Wald formula]
\label{P10:rem:acoustic_wald}
The acoustic metric $ds^2_\mathrm{ac} = \rho_0\,ds^2_\mathrm{Kerr}$
(Paper~VII~\cite{Paper7}) governs sound propagation and fixes the
Hawking temperature $T_H$ via the acoustic surface gravity $\kappa_H$.
It does \emph{not} enter the Wald entropy calculation, which uses
the \emph{physical} (Kerr) metric $g_{ab}$ in which the action
$S_\mathrm{tot}$ is written.
The conformal factor $\rho_0$ appears only in the propagator for
$\rho$-fluctuations, not in $\partial\mathcal{L}/\partial R_{abcd}$.
The two results therefore combine cleanly:
\begin{equation}
  T_H \;=\; \frac{\hbar\kappa_H}{2\pi k_B}
  \quad\text{(Paper~VII)},
  \qquad
  S \;=\; \frac{A_H}{4G_N}
  \quad\text{(Theorem~\ref{P10:thm:wald})},
  \label{P10:eq:TS_combined}
\end{equation}
reproducing the standard Bekenstein--Hawking thermodynamics exactly
within the condensate framework.
\end{remark}
\status{published}
\source{P10:rem:acoustic_wald}

\begin{remark}[Scope of the result]
\label{P10:rem:wald_scope}
Theorem~\ref{P10:thm:wald} holds for the full non-linear $k$-essence
Lagrangian~\eqref{P10:eq:parent_action_wald}: the argument in Step~2 requires only that
$\mathcal{L}_\mathrm{cond}$ is a function of $X$ alone, without
explicit Riemann-tensor dependence.
Higher-order curvature corrections (e.g.\ Gauss--Bonnet terms from
integrating out massive modes) would modify the result; no such terms
appear at the level of the condensate action of Papers~I--VII.
\end{remark}
\status{published}
\source{P10:rem:wald_scope}

\begin{lemma}[Conjugate momentum]
\label{P10:lem:momentum}
The canonical momentum conjugate to $\rho(x,t)$ from~\eqref{P10:eq:L_full} is
\begin{equation}
  \pi(x,t)
  \;=\; \frac{\partial\mathcal{L}}{\partial\dot\rho}
  \;=\; \frac{2\dot\rho}{\vp^6\bigl(1+\bst(\dot\rho^2-|\nabla\rho|^2)\bigr)^2}.
  \label{P10:eq:momentum}
\end{equation}
In the static limit $\dot\rho\to0$:
$\pi(x)\to 0$, and the linearisation about the saddle $f_*$ gives
\begin{equation}
  \pi(x) \;\approx\; \frac{2}{\vp^6(1+\bst\,|\nabla f_*|^2)^2}\,\dot\rho(x)
  \;\equiv\; W(x)\,\dot\rho(x),
  \label{P10:eq:momentum_linear}
\end{equation}
where $W(x) > 0$ is the \emph{symplectic weight function} at the saddle.
\end{lemma}
\status{published}
\source{P10:lem:momentum}

\begin{lemma}[Hamiltonian vector field of $\Efb$]
\label{P10:lem:ham_vf}
The Hamiltonian vector field of $\Efb$ on $(\CQr,\bOmega^{\mathrm{red}})$
is the condensate gradient flow with the scale mode projected out:
\begin{equation}
  X_E \;=\; -W^{-1}\left(\frac{\delta\Efb}{\delta f}
  - \frac{\langle\delta E/\delta f,\,g\rangle_W}
         {\langle g,g\rangle_W}\cdot g\right),
  \label{P10:eq:ham_vf}
\end{equation}
where $\langle u,v\rangle_W = \int W(x)\,u(x)v(x)\,d^3x$
and $g = -(r\partial_r f + z\partial_z f)$ is the Derrick mode.
\end{lemma}
\status{published}
\source{P10:lem:ham_vf}

\begin{remark}
\label{P10:rem:ham_vf_confirms}
Lemma~\ref{P10:lem:ham_vf} confirms that the scale-mode force projection
of Paper~IX is the correct gradient flow on the reduced phase space:
it is the Hamiltonian vector field of $\Efb$ on $\CQr$,
not an ad hoc numerical trick.
\end{remark}
\status{published}
\source{P10:rem:ham_vf_confirms}

\begin{lemma}[$q$-numbers are even Fibonacci numbers]
\label{P10:lem:fibonacci}
With $q = \vp^2$,
\begin{equation}
  \boxed{[n]_{\vp^2} \;=\; F_{2n}}
  \label{P10:eq:fibonacci}
\end{equation}
for all $n\geq0$, where $F_k$ denotes the $k$-th Fibonacci number
$(F_0=0,\,F_1=1,\,F_k = F_{k-1}+F_{k-2})$.
\end{lemma}
\status{published}
\source{P10:lem:fibonacci}

\begin{remark}
\label{P10:rem:qnumber_identity}
The identity $\vp^2 - \vp^{-2} = \sqrt{5}$ follows from
$\vp^2 = \vp+1$ and $\vp^{-2} = 2-\vp$, giving
$\vp^2-\vp^{-2} = (2\vp-1) = \sqrt{5}$.
The first few $q$-numbers are $[0]=0$, $[1]=1$, $[2]=3$, $[3]=8$,
$[4]=21$, $[5]=55$, $[6]=144$ --- every other Fibonacci number.
\end{remark}
\status{published}
\source{P10:rem:qnumber_identity}

\begin{corollary}[Well-definedness of the Liouville measure]
\label{P10:cor:measure_welldefined}
The regularised Liouville measure
\begin{equation}
  \boxed{\dmu = \det(\mathcal{L}_*)^{-1/2}\prod_{k\geq0}d\phi_k
  = \vp^{1/12}\,E_0^{-1/4}\,\mathcal{D}\phi}
  \label{P10:eq:measure_normalised}
\end{equation}
is well-defined and finite.
In particular, the Hilbert space $\Hil = L^2(\CQr,\dmu)$ of
Theorem~\ref{P10:thm:quantisation} has a canonically normalised
inner product.
\end{corollary}
\status{published}
\source{P10:cor:measure_welldefined}

\begin{remark}[The golden-ratio functional determinant]
\label{P10:rem:golden_functional_det}
Equation~\eqref{P10:eq:det_E01} is a remarkable result: the
\emph{zeta-regularised functional determinant of the Hopfion
stability operator is an exact power of the golden ratio},
$\det(\mathcal{L}_*) = \vp^{-1/6}$.
This could not have been anticipated from general principles;
it is a consequence of the golden-ratio spectrum $\lambda_k = E_0\vp^{2k}$
feeding into the geometric-series zeta function~\eqref{P10:eq:zeta_exact},
whose derivative at $s=0$ is controlled by the Bernoulli coefficient
$B_2 = 1/6$ at order $z^2/12 = \log(\vp)\cdot s/6$.
The exponent $-1/6$ is the second Bernoulli number
($B_2 = 1/6$) multiplied by $-1$: $\det(\mathcal{L}_*) = \vp^{-B_2}$.
\end{remark}
\status{published}
\source{P10:rem:golden_functional_det}

\begin{corollary}[Closed-form transition amplitudes]
\label{P10:cor:amplitudes}
The amplitude for a $k$-step descent $|n\rangle\to|n-k\rangle$ is
\begin{equation}
  \boxed{\langle n-k|\,a^k\,|n\rangle
  \;=\; \sqrt{\frac{[n]_q!}{[n-k]_q!}}
  \;=\; \sqrt{F_{2n}\cdot F_{2n-2}\cdots F_{2(n-k+1)}}},
  \label{P10:eq:amplitude}
\end{equation}
where the product uses $[j]_q = F_{2j}$ (Lemma~\ref{P10:lem:fibonacci}).
In particular, the amplitude for complete de-excitation from level $n$
to the ground state is
\begin{equation}
  \langle 0|\,a^n\,|n\rangle
  \;=\; \sqrt{[n]_q!}
  \;=\; \sqrt{F_2\cdot F_4\cdots F_{2n}}
  \;=\; 1,\;1,\;\sqrt{3},\;\sqrt{24},\;\sqrt{504},\;\ldots
  \label{P10:eq:full_deexcitation}
\end{equation}
and the Born-rule transition probability is
\begin{equation}
  P(n\to 0) = |\langle 0|a^n|n\rangle|^2
  = [n]_q! = \prod_{k=1}^n F_{2k}.
  \label{P10:eq:prob_n0}
\end{equation}
\end{corollary}
\status{published}
\source{P10:cor:amplitudes}

\begin{corollary}[$q$-coherent states and the generating function]
\label{P10:cor:coherent}
The $q$-coherent state $|z\rangle_q$ satisfying $a|z\rangle_q = z|z\rangle_q$ has
expansion
\begin{equation}
  |z\rangle_q
  \;=\; \mathcal{N}(|z|)\sum_{n=0}^\infty
  \frac{z^n}{\sqrt{[n]_q!}}\,|n\rangle
  \;=\; \mathcal{N}(|z|)\sum_{n=0}^\infty
  \frac{z^n}{\sqrt{F_2\,F_4\cdots F_{2n}}}\,|n\rangle,
  \label{P10:eq:q_coherent}
\end{equation}
where $\mathcal{N}(|z|) = \exp_q(-|z|^2/2)^{-1/2}$ is the
$q$-exponential normalisation and
$\exp_q(x) = \sum_{n=0}^\infty x^n/[n]_q!$
is the $q$-exponential generating function for Fibonacci ratios.
The overlap between two $q$-coherent states is
\begin{equation}
  {}_q\langle w|z\rangle_q
  \;=\; \frac{\exp_q(\bar{w}z)}{\sqrt{\exp_q(|w|^2)\,\exp_q(|z|^2)}}.
  \label{P10:eq:coherent_overlap}
\end{equation}
\end{corollary}
\status{published}
\source{P10:cor:coherent}

\begin{definition}[Configuration space]
\label{P10:def:config}
The \emph{condensate configuration space} is
\begin{equation}
  \mathcal{C}
  \;\equiv\; \left\{f:\RR^3\to[0,\pi]
  \;\Big|\;
  f\big|_{r=0}=\pi,\quad
  f\big|_{r,z\to\infty}=0,\quad
  \|f\|_{H^1}<\infty
  \right\},
  \label{P10:eq:config}
\end{equation}
equipped with the $H^1$ Sobolev norm.
The \emph{Hopf charge} of $f\in\mathcal{C}$ is the integer
\begin{equation}
  Q[f] \;\equiv\; \frac{1}{2\pi^2}\int_{\RR^3}
  \hat{n}\cdot(\partial_r\hat{n}\times\partial_z\hat{n})\,
  r\,dr\,dz\,d\Phi
  \;\in\; \ZZ,
  \label{P10:eq:hopf_charge}
\end{equation}
and the \emph{$Q$-sector} is $\CQ \equiv \{f\in\mathcal{C} : Q[f]=Q\}$.
\end{definition}
\status{published}
\source{P10:def:config}

\begin{definition}[Topological U(1) action]
\label{P10:def:U1}
The \textsc{u}(1) group acts on $\CQr$ by rotating the
winding phase:
\begin{equation}
  e^{i\theta}\cdot [f]
  \;=\; [f_\theta],
  \qquad
  f_\theta(r,z) = f(r,z),\quad
  \Phi_\theta = \Phi + Q\theta/Q = \Phi + \theta,
  \label{P10:eq:U1_action}
\end{equation}
where $\Phi$ is the Hopf azimuthal phase and $Q\in\ZZ$ is the
Hopf charge~\cite{Paper1,Paper9}.
The moment map of this action is
\begin{equation}
  \mu: \CQr \to \RR,
  \qquad
  \mu([f]) = Q\cdot\omega([f]),
  \label{P10:eq:moment}
\end{equation}
where $\omega([f])\in[0,2\pi]$ is the solid angle subtended by
the Hopfion configuration at the observer's position.
\end{definition}
\status{published}
\source{P10:def:U1}

\begin{definition}[$q$-number and $q$-factorial]
\label{P10:def:q_numbers}
For $q = \vp^2$ and $n\in\ZZ_{\geq0}$, define the $q$-number
\begin{equation}
  [n]_q \;\equiv\; \frac{q^n - q^{-n}}{q - q^{-1}},
  \qquad
  [n]_q! \;\equiv\; \prod_{k=1}^n [k]_q,
  \qquad [0]_q! = 1.
  \label{P10:eq:q_numbers}
\end{equation}
\end{definition}
\status{published}
\source{P10:def:q_numbers}

\begin{equation}
  \Gamma_\mathrm{1-loop}
  \;=\; \tfrac{1}{2}\log\det\mathcal{L}_*
  \;=\; -\tfrac{1}{12}\log\vp,
  \label{P10:eq:Gamma_1loop}
\end{equation}

\begin{equation}
  \hat{H}_\mathrm{fb}
  \;\equiv\; -i\hbar\nabla_{X_E} + \Efb,
  \label{P10:eq:Hfb}
\end{equation}

\begin{equation}
  \mathcal{L}[\rho]
  \;=\; \frac{\dot\rho^2 - |\nabla\rho|^2}{\vp^6
         \bigl(1+\bst(\dot\rho^2-|\nabla\rho|^2)\bigr)}.
  \label{P10:eq:L_full}
\end{equation}

\begin{equation}
  \mathcal{L}_\mathrm{kin}[\rho]
  \;=\; \frac{|\nabla\rho|^2}{\vp^6\bigl(1+\bst\,|\nabla\rho|^2\bigr)},
  \label{P10:eq:Lkin}
\end{equation}

\begin{equation}
  E_k = E_0\cdot\exp\!\bigl(2\log\vp\cdot k\bigr) = E_0\,\vp^{2k},
  \label{P10:eq:WZW_energy}
\end{equation}

\begin{equation}
  S_\mathrm{tot}[g,\rho]
  \;=\; \underbrace{\frac{1}{16\pi G_N}\int\!\sqrt{-g}\;R\;\mathrm{d}^4x}_{S_\mathrm{EH}}
  \;+\; \underbrace{\int\!\sqrt{-g}\;P(X)\;\mathrm{d}^4x}_{S_\mathrm{cond}},
  \label{P10:eq:action_total}
\end{equation}

\begin{align}
  F_{2n}
  &= \frac{\vp^{2n} - \psi^{2n}}{\sqrt{5}}
   = \frac{\vp^{2n} - \vp^{-2n}}{\sqrt{5}}.
  \label{P10:eq:binet_even}
\end{align}

\begin{equation}
  \frac{\partial\mathcal{L}_\mathrm{EH}}{\partial R_{abcd}}
  \;=\; -\frac{1}{32\pi G_N}
  \bigl(g^{ac}g^{bd} - g^{ad}g^{bc}\bigr).
  \label{P10:eq:dLEH}
\end{equation}

\begin{equation}
  \frac{\partial\mathcal{L}_\mathrm{cond}}{\partial R_{abcd}}
  \;=\; P'(X)\,\frac{\partial X}{\partial R_{abcd}}
  \;=\; 0,
  \label{P10:eq:dLcond}
\end{equation}

\begin{equation}
  \log\det(\mathcal{L}_*)\big|_{E_0=1}
  = -\frac{\log\vp}{6},
  \qquad
  \boxed{\det(\mathcal{L}_*) = \vp^{-1/6}}.
  \label{P10:eq:det_E01}
\end{equation}

\begin{equation}
  \prod_{k=0}^\infty \lambda_k^{-1/2}
  = E_0^{-N/2}\,\prod_{k=0}^\infty\vp^{-k}
  = E_0^{-N/2}\,\vp^{-N(N-1)/2}\;\to\; \vp^{-1/12}/\eta(\tau),
  \label{P10:eq:eta_connection}
\end{equation}

\begin{proposition}[Logarithmic heat kernel]
\label{P10:prop:heat_kernel}
The heat kernel
$\Theta(t) = \mathrm{Tr}[e^{-t\mathcal{L}_*}]
= \sum_{n=0}^\infty e^{-\vp^{2n}E_0 t}$
satisfies
\begin{equation}
  \Theta(t)
  \;\sim\; -\frac{\log(E_0 t)+\gamma}{2\log\vp}
  \qquad\text{as }t\to0^+,
  \label{P10:eq:heat_kernel}
\end{equation}
where $\gamma = 0.5772\ldots$ is the Euler--Mascheroni constant.
This is a $\log(1/t)$ divergence, parametrically softer than the
$t^{-3/2}$ divergence of a standard scalar field in three dimensions.
\end{proposition}
\status{published}
\source{P10:prop:heat_kernel}

\begin{remark}[UV softness and the spectral zeta residue]
\label{P10:rem:uv_softness_residue}
The coefficient $1/(2\log\vp)$ in~\eqref{P10:eq:heat_kernel} equals
$\mathrm{Res}_{s=0}\,\zeta_{\mathcal{L}_*}(s)$
(Theorem~\ref{P10:thm:spectral_zeta}).
The $\log(1/t)$ behaviour (one log-divergent heat-kernel coefficient)
replaces the three power-law-divergent coefficients $a_0,a_{1/2},a_1$
of standard QFT: the golden-ratio discretisation
makes the spectrum too sparse for the Weyl law to hold, softening
the UV behaviour by a power law.
The $B_2=1/6$ factor in $\det\mathcal{L}_*=\vp^{-B_2}$ and the
$1/(2\log\vp)$ factor in~\eqref{P10:eq:heat_kernel} are two manifestations
of the same Bernoulli-series regularisation.
\end{remark}
\status{published}
\source{P10:rem:uv_softness_residue}

\begin{equation}
  \frac{1}{2\pi}\oint_\gamma \bOmega^{\mathrm{red}}
  \;=\; \frac{Q}{2\pi}\oint_\gamma d\mu
  \;=\; Q\cdot 1
  \;=\; Q \;\in\;\ZZ,
  \label{P10:eq:integrality}
\end{equation}

\begin{equation}
  \log\det(\mathcal{L}_*)
  = -\zeta'_{\mathcal{L}_*}(0)
  = -\frac{\log\vp}{6} + \frac{\log E_0}{2}
    - \frac{(\log E_0)^2}{4\log\vp}.
  \label{P10:eq:logdet}
\end{equation}

\begin{equation}
  P(X) \;=\; \frac{X}{\vp^6\,(1+\bst X)}
  \label{P10:eq:parent_action_wald}
\end{equation}

\begin{align}
  [n+1]_q - q[n]_q
  &= \frac{q^{n+1}-q^{-(n+1)}-q(q^n-q^{-n})}{q-q^{-1}} \notag\\
  &= \frac{q^{n+1}-q^{-(n+1)}-q^{n+1}+q^{1-n}}{q-q^{-1}} \notag\\
  &= \frac{q^{1-n}-q^{-(n+1)}}{q-q^{-1}}
   = \frac{q^{-n}(q-q^{-1})}{q-q^{-1}}
   = q^{-n}. \label{P10:eq:q_identity}
\end{align}

\begin{equation}
  \dmu \;=\; \prod_k \lambda_k^{-1/2}\,d\phi_k,
  \label{P10:eq:regularised_measure}
\end{equation}

\begin{equation}
  \left.\frac{d\Efb}{dt}\right|_{f=f_*} = 0,
  \qquad
  \sopt[f_*] \;=\; 2^{1/6},
  \label{P10:eq:saddle_conditions}
\end{equation}

\begin{equation}
  \delta^2\Efb[\phi]
  = \int \phi\,\mathcal{L}_*[\phi]\,d^3x,
  \qquad
  \mathcal{L}_* = -\Delta + V_*,
  \label{P10:eq:second_var}
\end{equation}

\begin{equation}
  \mathcal{L}_*
  \;\equiv\; -\Delta + V_*(x),
  \qquad
  V_*(x) = -\frac{d^2\Efb}{df^2}\bigg|_{f=f_*},
  \label{P10:eq:stab_op}
\end{equation}

\begin{equation}
  \Theta(t)\;\approx\;
  \frac{1}{2\log\vp}\int_{E_0 t}^\infty\frac{e^{-u}}{u}\,du
  \;=\; \frac{\Gamma(0,\,E_0 t)}{2\log\vp},
  \label{P10:eq:theta_gamma}
\end{equation}

\begin{equation}
  S_\mathrm{Wald}
  \;=\; -2\pi\int_{\mathcal{H}}
  \frac{\partial\mathcal{L}}{\partial R_{abcd}}\,
  \varepsilon_{ab}\,\varepsilon_{cd}\;\mathrm{d}^2x,
  \label{P10:eq:wald_formula}
\end{equation}

\begin{align}
  \zeta_{\mathcal{L}_*}(s)
  &= \frac{1}{2\log\vp\cdot s}
  + \underbrace{\!\left[\frac{1}{2} - \frac{\log E_0}{2\log\vp}\right]}_{=\,\zeta_\mathrm{reg}(0)}
  + \underbrace{\!\left[\frac{\log\vp}{6} - \frac{\log E_0}{2}
    + \frac{(\log E_0)^2}{4\log\vp}\right]}_{=\,\zeta'(0)}\!\cdot s
  + O(s^2).
  \label{P10:eq:zeta_coeffs}
\end{align}

\begin{equation}
  \zeta_{\mathcal{L}_*}(s)
  = \sum_{k=0}^{\infty}(E_0\vp^{2k})^{-s}
  = E_0^{-s}\sum_{k=0}^{\infty}\vp^{-2ks}
  = \frac{E_0^{-s}}{1-\vp^{-2s}},
  \label{P10:eq:zeta_derivation}
\end{equation}

\begin{align}
  \zeta_{\mathcal{L}_*}(s)
  &= E_0^{-s}\!\left[\frac{1}{2\log\vp\cdot s}
   + \frac{1}{2} + \frac{\log\vp\cdot s}{6} + O(s^2)\right].
  \label{P10:eq:zeta_laurent}
\end{align}

